% 規劃文件；尚未執行，不可作為第四章的新實驗結果。
% LaTeX 片段，使用既有論文導言中的 graphicx、hyperref 與數學套件。
% 本文件不載入既有第四、第五章，也不修改已寫入的實驗數據。

\section*{純模擬實驗重建規劃}

本規劃以重新蒐集資料為前提，既有紀錄僅供開發、估計變異及檢查流程。
核心版預計執行 $949$ 趟，完整版為 $4302$ 趟；
另需 $300$ 或 $600$ 張獨立測試影像。
兩個數字是本研究範圍下的工作預算，並非論文審查或統計顯著性的最低門檻。
工程試跑、基線調參與視覺訓練資料另計。
核心版的 $24$ 趟模型資料包含訓練、驗證與測試，不能全部視為獨立測試樣本。

\subsection*{一、需要證明的主張}

\begin{enumerate}
    \item 執行模型能改善對位置命令介面的預測，且在線上操作中具有可辨識的效果。
    \item 動態障礙預測、CBF、原始掃描防護及參考調整，對到達、淨空及時間有不同作用。
    \item 全身預測控制相較外部反應式全身控制，在共同任務與限制下具有可量測的優點與代價。
    \item 移動與操作重疊能在完成條件相同的情況下改善任務效率。
    \item 視覺估測、命令交接與機構操作能於完整住宅任務中接續，並可說明失敗來源。
\end{enumerate}

底盤 SE(2) GMPC 是既有方法基礎。
若不將幾何表示本身列為新貢獻，無須先重做大量幾何追蹤優越性實驗；
其獨立追蹤比較列為完整版的 N0。
各方法的優越性須由結果判定，不能將預期收益預先寫成結論。

\subsection*{二、相關研究實際採用的驗證方式}

\begin{enumerate}
    \item Tang 等人之 GMPC：比較 NMPC 與回授線性化控制，
    使用圓形與蝴蝶軌跡、隨機起始位姿及兩種輪式平台，另比較線性化方式。
    參考：\url{https://arxiv.org/html/2403.07317v1}，第 IV 節。
    本次核對的實驗正文未明列 Monte Carlo 的確切趟數，故不推估其總數。

    \item BARN：$300$ 個模擬環境，每個環境各執行 DWA、E-Band 五次，
    共 $3000$ 趟模擬導航；另有 $50$ 趟實體測試。
    其啟示是增加環境多樣性，並保留失敗懲罰與路長正規化。
    參考：\url{https://www.cs.utexas.edu/~xiao/papers/navdiffdataset_ssrr.pdf}，第 IV 節。

    \item DynaBARN：建立 $300$ 個動態環境，測試障礙速度、速度變化與路徑形狀。
    比較 DWA、TEB、TD3 與行為複製；各基線於每個受測環境部署十次。
    不能直接以資料庫總數乘十，推定正文所有基線實際測試總量。
    參考：\url{https://people.cs.gmu.edu/~xiao/papers/dynabarn.pdf}，第 III--IV 節。

    \item Haviland 等人之 Holistic Mobile Manipulation：
    使用 $1000$ 個隨機可達末端目標，比較整體式控制與分離規劃，並做參數消融。
    實體取放為十輪、合計 $100$ 個物件，包含 $113$ 次抓取嘗試。
    任務完成數與首次抓取成功數應分別統計。
    參考：\url{https://arxiv.org/pdf/2109.04749}，第 V--VI 節。

    \item Burgess-Limerick 等人之 MotM 架構：
    取放測試每個方法與速度配置重複十次，區分靜態與移動物件，
    比較序列反應式與規劃式基線，並分析時間、成功及加速度。
    提出方法的兩項實體實驗合計 $120$ 次嘗試，不能解讀為所有基線合計 $120$ 次。
    參考：\url{https://arxiv.org/pdf/2212.06991}，第 IV--V 節。

    \item Reactive Base Control for MotM：
    使用 $50$ 組隨機物件排列進行模擬，並將相同排列加入兩個靜態障礙後重測。
    動態障礙部分另有案例與距離分析，不能據此套用同樣的配對時間結論。
    參考：\url{https://arxiv.org/pdf/2309.09393}，第 VI--VII 節。

    \item Mittal 等人之關節式物件操作：
    研究不同廚房配置、抽屜及可動家具，分別處理物件計畫與全身控制，
    並評估靜態及動態場景。本次僅據作者專案確認驗證範圍，不填入未核對的趟數。
    參考：\url{https://www.pair.toronto.edu/articulated-mm/}。
\end{enumerate}

\subsection*{三、基線與公平比較}

導航主比較採四組：本研究完整導航、MPPI Omni、DWB Omni 與 RPP。
DWB 是新增基線的優先選項，因 Nav2 Jazzy 的官方實作支援全向運動。
RPP 保留為常見的前進與轉向參考，其側向命令為零，不能解釋為同運動能力的純演算法比較。
官方說明：\url{https://api.nav2.org/nav2-jazzy/html/md_nav2_dwb_controller_README.html}。
TEB 可另外研究，但本計畫不以移植另一套控制器取代必要的模型與整機驗證。

導航四組共享感測、外形、全域路徑來源、進度與到達條件、
執行端平滑器及獨立原始掃描防護。
MPPI、DWB 與 RPP 的表格名稱須註明共用防護，不能將防護後結果稱為原生控制器表現。
消融 N2 才獨立改變防護與 CBF，並以模擬接觸資料量測結果。
共同防護的介入比例、介入時長及介入前後命令差也列為指標。

全向方法共享合成平移速度上限
$\sqrt{(v_x^B)^2+(v_y^B)^2}\leq\bar v$，
並共享輪速、輪級命令變化及偏航上限，避免逐軸限速形成對角線速度優勢。
各基線可在相同開發資料及調參預算內使用適合自己的內部成本，正式測試後不再調整。
評估器只讀取事後真值，不提供額外障礙或目標資訊給任一主比較方法。

局部操作主比較採本研究 WB-MPC、移植的 Holistic 反應式全身 QP，
以及嚴格停駐策略。
反應式 QP 依作者公開方法移植至本平台，保留共同手臂位置介面及命令上限；
需列出速度命令轉位置設定點的轉接方式，並先檢查其自由空間響應。
這是外部方法的移植對照，不能宣稱直接重現原論文數字。
公開實作：\url{https://jhavl.github.io/holistic/}。
本研究的 $N=1$ 僅為內部時域消融，不能替代此一外部基線。

局部嚴格停駐比較使用整段機構運動皆有固定底盤可行解的共同場景，
以事前幾何檢查定義，不依實驗成敗事後刪樣。
住宅序列基線則在接近前停車，後續機構操作允許全身協同；
兩種停駐條件必須分別命名。

\subsection*{四、核心版的執行矩陣}

下表均為新增執行數，已扣除能完全共用的完整方法參考趟數。
每個種子代表預先保存的完整初始狀態、路線及擾動案例，
不是將相同確定性執行任意重播十次後視為十個獨立場景。

\begin{table}[htbp]
    \centering
    \caption{核心版實驗預算}
    \small
    \renewcommand{\arraystretch}{1.2}
    \begin{tabular}{l p{0.65\textwidth} r}
        \hline
        編號 & 設計 & 新增趟數 \\
        \hline
        M0 & 兩種負載、三個構型區域、四段獨立激勵紀錄 & $24$ \\
        N1 & 三類布局、兩類障礙條件、十個種子、四種導航方法 & $240$ \\
        N2 & 兩個關鍵交會場景、十個種子、六個新增消融配置 & $120$ \\
        W1 & 四種操作、兩種空間條件、十個種子、三種策略 & $240$ \\
        W2 & 圓柱與抽屜、各十個種子、八個新增消融配置 & $160$ \\
        F0 & 三個階段、三種故障、五個種子 & $45$ \\
        E0 & 兩個住宅布局、兩類行人條件、十個種子、三種策略 & $120$ \\
        \hline
        合計 & 另有 $300$ 張測試影像；開發試跑及調參另計 & $949$ \\
        \hline
    \end{tabular}
\end{table}

M0 每一激勵紀錄均包含階躍、斜坡、多頻激勵及多關節運動。
六個負載與構型分層中，各保留一段訓練、一段驗證及兩段測試，
合計為六段訓練、六段驗證及十二段測試。
以完整紀錄分割，不能將相鄰時間窗隨機分到訓練與測試。
測試 $50$、$100$、$250$ 與 $500$ ms 預測，
分開呈現使用後續實錄輸入的條件重播，以及只依當時可用計畫推進的預測。
線上偏置僅能利用預測起點以前的已完成區間更新。

N1 的三類布局為開闊通行、狹窄轉角與家具雜亂區，
各包含靜態障礙及動態交會。
在案例層級固定路線長度範圍與最小通道裕度，
動態種子涵蓋橫越、迎面或從側方匯入，並保存起跑時的相位。
正式案例由另外的開發布局與種子產生，不從測試集合挑選較容易的路線。

N2 共比較七個配置，完整配置的結果取自相同 N1 案例，故只增加六組：
\begin{enumerate}
    \item GMPC，關閉 CBF、防護及參考調整。
    \item GMPC 加預測式 CBF，關閉防護及參考調整。
    \item GMPC 加原始掃描防護，關閉 CBF 及參考調整。
    \item GMPC 加預測式 CBF 與防護，關閉參考調整。
    \item 第四組改用零障礙速度，保留同樣的 CBF 形式。
    \item 第四組加入速度調整，關閉側移候選。
    \item 完整方法：速度與側移調整、預測式 CBF 及防護。
\end{enumerate}
前四組提供 CBF 與防護的二乘二消融，
第四與第五組檢查障礙運動估測，第六與第七組檢查側移的增量效果。
四個方法不得使用不同的基礎間距與命令上限來製造優勢。

W1 的四種操作為圓柱取放、筆狀物取放、門扇開關及抽屜開關。
兩類空間為一般空間與家具附近的受限空間。
局部門扇先採 $45^\circ$ 開啟目標，以建立嚴格停駐仍可完成的共同工作空間；
住宅任務另採 $85^\circ$，其序列策略允許機構操作期間的底盤配合。
每趟包含共同進場、接近、操作及退出事件；
取放另包含抬升、運送、釋放與收臂。
物件尺寸、位置及朝向在測試範圍內變動，並為所有方法保存相同初始狀態。

W2 的完整方法加上八個新增配置為：直接設定點變化率積分模型、
關閉在途命令前推、固定偏置、關閉末端積分、兩種補償皆關閉，
以及三個其餘預測步數。
名目預測步數先由開發集在 $\{1,5,10,20\}$ 中選定，
正式比較時固定成本定義、每步時間尺度、終端成本矩陣與約束。
各時域均由同一參考生成器提供對應未來時刻的目標。
成本累加方式與終端時刻仍隨時域改變，故結果解讀為配置比較。
對有條件啟用的補償，另報告實際啟用時間及有效樣本數。

F0 在導航、接近及抽屜拉開三個階段，
分別注入工作程序逾時、QP 不可行及感知資料逾時。
注入後持續物理步進，核對底盤及關節的實際速度，
不以停止尾段事件或模擬凍結代替停穩。
記錄故障、最後非零命令、輸入排清及實際停穩四個時刻。

E0 為完整開門、通行、圓柱與筆取放收納、抽屜關閉、回程關門及返回。
第二個住宅布局不參與參數調整；
行人採事前保存的外生軌跡，不因方法不同而另改讓行規則。
若另研究會主動讓行的行人模型，須另列條件。

\subsection*{五、完整版的執行矩陣}

完整版包含核心版案例，不能將兩者的總數相加。
僅在控制版本、資產、判準與資料分割均未變更時，才可直接沿用核心版紀錄。

\begin{table}[htbp]
    \centering
    \caption{完整版實驗預算，均為去重後的執行數}
    \small
    \renewcommand{\arraystretch}{1.15}
    \begin{tabular}{l p{0.65\textwidth} r}
        \hline
        編號 & 設計 & 趟數 \\
        \hline
        M0 & 兩種負載、三個構型區域、十二段激勵 & $72$ \\
        N0 & 四種軌跡、三級初始偏差、十個種子、三控制器 & $360$ \\
        N1 & 四類布局、三類障礙條件、二十種子、四方法 & $960$ \\
        N2 & 四個關鍵場景、二十種子、六個新增配置 & $480$ \\
        W0 & 四類自由空間參考、二十種子、兩控制器 & $160$ \\
        W1 & 四操作、三類空間與動態條件、二十種子、三策略 & $720$ \\
        W2 & 四操作、二十種子、八個新增消融配置 & $640$ \\
        V1 & 四操作、二十種子，改用真值目標作診斷 & $80$ \\
        C1 & 四操作、二十種子，改用固定相對位姿接觸對照 & $80$ \\
        R0 & 四操作、二十種子、四個新增擾動條件 & $320$ \\
        S1 & 四操作、二十種子，改用最近可行底盤選位 & $80$ \\
        K1 & 門扇操作二十種子，移除位姿容許範圍限制 & $20$ \\
        F0 & 三階段、三故障、十種子 & $90$ \\
        E0 & 兩住宅布局、兩行人條件、二十種子、三策略 & $240$ \\
        \hline
        合計 & 另有 $600$ 張測試影像 & $4302$ \\
        \hline
    \end{tabular}
\end{table}

N0 比較 GMPC、具角度包絡處理的歐氏 MPC 與本體座標 PID，
軌跡包含圓形、八字、S 形與側移合併偏航。
W0 包含底盤平移時保持世界 TCP、底盤偏航時保持世界 TCP、
世界 TCP 路徑追蹤與底盤相對持物參考。

完整版 N1 加入第四個未見住宅布局及多障礙運動條件，
後者包含不同速度、轉彎及速度變化，不只增加障礙數量。
W1 的第三條件為動態障礙交會。
M0 的擴充保留核心版測試資料的角色，
每個負載與構型分層增至三段訓練、三段驗證與六段測試。

V1 僅替換操作目標資訊，與原視覺條件逐案配對，
用於估計感知對控制的影響，不能列為主方法的自主能力。
C1 比較接觸建模假設，固定相對位姿結果不與摩擦接觸混合統計。
R0 與名目條件配對，增加命令延遲至 $160$ ms、目標位置偏置 $20$ mm、
視覺中斷 $500$ ms 及手臂驅動增益乘 $0.8$ 四個條件，
各次只改一項；驅動增益乘數不直接稱為同百分比的響應速度變化。
這些數值是開發前的初始提案，須於開發階段確認合理性後固定。
S1 的最近可行基線保留相同導航、IK 與淨空可行性檢查，只替換排序準則。
K1 保留機構目標與接觸設定，只移除對應的硬位姿容許範圍。
若未宣稱選位成本或位姿容許範圍具有獨立效益，S1、K1 可不納入核心範圍。

\subsection*{六、接觸、視覺與計時條件}

主比較採物理摩擦抓持，使用可動物件、被動鉸鏈與滑軌，
由夾爪及物件接觸形成抓持與機構運動。
控制器仍可使用固定相對位姿作為幾何參考，
但主要評估的模擬器不以新增固定連接或直接更新物件位姿確保成功。
若尚未完成此接觸模型，先將結論限定為運動學協調，
不將固定相對位姿條件下的完成數稱為摩擦抓持可靠性。

視覺獨立測試以三十或六十個未見場景種子，各取十個相機視角，
得到三百或六百張 RGB-D 影像。
每類至少覆蓋五十個有效物件實例；不足時增加不同場景，不能只複製影格。
測試場景與訓練、驗證分割以場景及序列為單位，避免相鄰影格洩漏。
離線比較預測遮罩與真值遮罩下的相同幾何估測，
統計各類 mask AP、IoU、漏檢與誤檢，以及三維位置、方向及門角誤差。
各類幾何方向須處理圓柱旋轉對稱與筆軸正反等價。

物理與控制步長建議先沿用 $0.01$ s 與 $0.05$ s，
經開發確認後固定。
所有比較方法共用同一模擬配速，不能為較慢方法單獨放慢物理時間。
若一倍速下無法滿足期限，可採固定共同配速評估控制效果，
並將一倍速即時性留為未驗證主張。
計算耗時採主機時間，任務運動採模擬時間，另保存實際 RTF。
本規劃不要求提供個人電腦型號。

\subsection*{七、完成條件、指標與統計}

導航成功初始建議為位置誤差小於 $0.15$ m、朝向誤差小於 $10^\circ$，
並低速保持 $0.5$ s，且整段沒有非任務指定接觸。
局部與住宅門扇開啟參考分別可設為 $45^\circ$ 與 $85^\circ$，
容許差 $5^\circ$，關閉為小於 $3^\circ$；
抽屜拉出可設為 $200\pm5$ mm，關閉為小於 $5$ mm。
收納檢查物件幾何位於盒內並穩定保持一秒，
返回則同時檢查底盤位姿、低速與手臂收納構型。
上述門檻需依機構公差於開發集定案，正式測試後不得變更。
整體成功要求所有子任務、最終條件與非指定接觸條件共同通過。

計時由共同起跑事件至實際完成條件成立，
包含方法自己的規劃、選位、等待及恢復過程。
共同的模型載入與固定預熱可在起跑前完成，但須對所有方法一致。
主導航控制器比較可固定全域路徑；端到端任務則納入實際規劃與重規劃成本。
方法失敗、逾時、保護中止與紀錄基礎設施失效分別標註。
測試不能因方法失敗而無紀錄補跑；基礎設施失效亦保留原紀錄與補跑關聯。

主指標為無非指定接觸的完成率與失敗類別。
另列最小全身或底盤淨空、實際接觸、TCP 世界目標誤差、
相對機構誤差、抓持相對位移、關節及設定點餘量、
實際放行命令變化、總任務及分段時間、求解與來回時間、逾時與停止結果。
幾何負淨空與模擬接觸分開量測。
底盤平移加速度先將本體速度轉到世界座標再差分；
設定點變化率的差分只稱命令變化，不冒稱實際關節加速度。

任務時間先呈現兩方法均完成的配對時間及其樣本數，
再以所有嘗試共同呈現成功率及預先定義的失敗懲罰分數。
例如成功時取 $1-T/T_{\max}$，未成功取零，
可作為輔助任務分數，但不能稱為實際完成時間。
各任務的 $T_{\max}$ 由開發集定案，不能依方法或測試結果調整。

時間序列、預測窗口與相鄰影格不視為獨立樣本。
統計以場景種子為配對或群聚單位，
對連續量採配對差、相對時間差及群聚 bootstrap 信賴區間，
對配對成功結果可使用 exact McNemar 檢定。
多個主要基線比較需預先指定主對照，並作 Holm 等多重比較調整。
各布局與任務分層呈現，另外使用事先約定權重計算整體結果。

趟數需先看研究問題與變異。
以獨立同分布二項試驗作示例，十趟皆成功時，
失敗率的單側 $95\%$ 上限仍約為 $25.9\%$；
六十趟零失敗時約為 $4.87\%$。
因此十趟足以作為一個配置的起始重複數，但不能據此保證高可靠性。
精確二項區間說明：
\url{https://itl.nist.gov/div898/handbook/prc/section2/prc241.htm}。
若目標是證明約 $5\%$ 的時間差，
而配對對數時間差標準差約為 $0.15$，
以常態近似的雙側 $5\%$ 檢定與 $80\%$ 檢定力估算需約七十對；
此為示例，應由先導資料重估，亦不能將成功選樣造成的偏差以增加趟數消除。

\subsection*{八、資料交付與執行順序}

每趟至少交付：版本及設定雜湊、資產及場景快照、初始狀態、
所有隨機種子、觀測時間與相機狀態、RGB-D 或可追溯影格、
底盤位姿、實際關節及設定點、提案與最終放行命令、
排程版本及生效步、障礙估計與事後真值、接觸資料、
任務事件、求解狀態、逾時及實際停止資料。
原始命令、感知輸入與評分真值的用途需分開記錄。
基本物理訊號至少依每個物理步保存，控制資料每次更新保存，影像保留實際觀測頻率。

執行順序為：
模型與紀錄介面核對、基線自由空間校核、M0 辨識與資料分割、
開發集調參、配置凍結、N1/N2、W1/W2、F0，最後 E0。
完整版再擴充 V1、C1、R0、S1、K1 與更多布局和種子。
同案例的完整方法與消融變體放在同一排程區塊，隨機化方法順序。
CSV 的 within\_case\_order 已提供可重現的區塊內順序；
此清單僅為規劃，資產及設定仍待凍結，不能直接視為可啟動的批次工具。

第一批建議為十二趟先導：四方法各一趟導航、
三策略各一趟圓柱及抽屜操作，共六趟，另加兩趟故障停止。
這批用於確認比較與資料完整性，不併入正式測試。
正式測試進行中若修正控制器，需建立新版本資料群，重新執行受影響的共同案例。
以先導量得的單趟主機耗時與重置時間估計排程，
不由目前尚未知的計算速度推定整批可在幾天內完成。

\subsection*{九、預計形成的論文圖表}

\begin{enumerate}
    \item 模型辨識參數及時間分離的預測誤差--時域曲線。
    \item 導航各場景的完成率、失敗分類、淨空分布與共同案例時間圖。
    \item CBF 與掃描防護消融，以及障礙預測、速度及側移的增量比較。
    \item 四種操作的基線比較、末端誤差與底盤手臂分工圖。
    \item 預測步數的追蹤、命令變化及計算成本曲線，及模型補償消融。
    \item 視覺各類別分割與三維估測誤差，完整版另列真值目標診斷。
    \item 故障至實際停穩的時間與速度曲線。
    \item 住宅完成率、階段失敗漏斗、所有嘗試分數與配對完成時間。
\end{enumerate}
